\chapter{理论基础与文献综述}

\section{概念界定}

\subsection{国际流动}

国际流动是一个涵盖范围广泛的概念，包括劳动力的跨国迁移、国际移民、难民流动、商务旅行以及学术人员的跨境流动等多种形式。在高等教育与科学研究领域，国际流动主要表现为国际学术流动，即科研人员、学者和学生以学术活动为目的的跨国流动。本文所讨论的国际流动特指国际学术流动。国际学术流动（International Academic Mobility）作为高等教育国际化的重要组成部分，近年来受到学界的广泛关注，学者们从不同视角对其概念进行了界定。Netz等人将国际学术流动界定为科学家及相关信息跨越国家边界的流动过程\cite{netz2020}，国际学术流动既可以表现为科研人员的物理流动，如出国学习、访学或工作，也可以表现为依托信息技术开展的虚拟学术交流与合作。Shen等人在此基础上进一步拓展了流动主体的范围，强调国际学术流动研究应将研究生，尤其是博士生纳入研究视野，并将国际学术流动界定为学生、研究人员与学者等主体的跨境移动，以及资源、知识与信息等要素的跨境流动\cite{shen2022}。在研究主体的界定上，李倩文等人将学术人员定义为以学术为志业的研究者，涵盖高校教师、科研机构研究人员、博士后群体以及在学术系统从事研究的博士研究生，相应地，国际学术流动被界定为学术人员以研究生升学、留学或职业变动的方式从某一国家机构转移到另一国家机构的过程\cite{liqianwen2023}。在流动形式的分类上，王艳和赵江南归纳了国际学术流动的主要类型\cite{wangyan2026}，包括学生出国攻读博士学位、开展研究访问、毕业生赴海外从事博士后研究以及教师出国研修等。张颖等人则进一步明确，在海外攻读学位、从事全职或非全职研究工作均属国际学术流动范畴\cite{zhangying2024}。综合上述界定，国际学术流动涵盖了科研人员跨越国界开展学习与研究活动的多种形式，既包括短期访学，也包括长期的求学或工作经历，流动主体亦从传统的教师群体延伸至博士生、博士后等早期职业研究者。

结合本研究的研究对象与研究目的，本文将国际学术流动定义为科研人员跨越国家边界进行的学术活动中的物理流动，具体指博士研究生在培养阶段赴海外高校或科研机构开展阶段性科研活动的跨国学术经历。

\subsection{学术资本}

马克思认为，资本不是一般意义上的物或财富，而是在特定历史条件下以物为载体表现出来的社会生产关系，其核心在于价值的不断增殖\cite{marx2004}。在这一意义上，资本的概念可用于理解不同社会场域中各类资源的积累、配置与转化。学术场域同样遵循这一逻辑，学者在学术活动中投入的资源在持续的学术实践中得到积累并转化为新的竞争优势，学术资本（Academic Capital）的概念由此而来。围绕这一概念，学者们从各自研究视角出发，作出了不同界定。Cuellar等人认为，学术资本是学者带给组织的能力与声誉资产的集合，体现为其在特定研究领域发展与传播学术思想的"资本库"\cite{cuellar2016}。Panofsky则从资本结构的角度指出，学术资本是科学专长、社会交往能力以及象征性与物质性资源的综合体，是科学家获得同行认可的重要基础\cite{panofsky2011}。国内一些学者基于学术场域的视角对这一概念进行了界定，胡钦晓认为，学术资本是在特定学术场域中，个体或组织凭借稀缺的专业知识与技能逐步形成的学术成就与学术声望，并能够在学术活动与资源交换过程中实现价值增值的资源总和\cite{huqinxiao2020}。张琳等人则将其概括为高校教师在学术场域中围绕科学知识生产与增长所积聚的各类学术资源\cite{zhanglin2024}。总体而言，学术资本是科研人员在学术场域中通过学术实践逐步积累的知识能力、社会关系与学术声誉等多维资源的综合体。

本研究在上述定义的基础上，结合研究对象的特殊性，将学术资本界定为CSC资助博士生在国际学术流动过程中积累的、有助于其学术成长与职业发展的多维资源总和，涵盖知识与技能、社会关系网络、声誉与认可以及心理特质等多个层面。关于本研究所采用的具体理论框架与各维度资本的界定，将在下一节理论基础中详细阐述。

\section{理论基础}

\subsection{布迪厄资本理论}

布迪厄（Pierre Bourdieu）资本理论源于其对社会不平等再生产机制的长期研究，是20世纪社会学中最具影响力的理论框架之一。与以往主要将资本理解为经济意义上物质资源的研究不同，布迪厄首次在社会理论层面系统地将"资本"概念加以扩展和概念化，将其引入文化、社会与象征等多个维度，从而为解释社会结构、权力关系及社会分层提供了新的分析工具。布迪厄在《资本的形式》（The Forms of Capital）中指出，资本并非静态的物或简单的资源存量，而是一种能够在特定社会条件下转化为优势和权力的社会力量\cite{bourdieu1986}。

在布迪厄看来，资本是"积累的劳动"，既可呈现为物质形态，也能以内化于个体自身的形式存在，并且只有在特定社会关系和制度环境中，资本才能发挥其效力。基于这一理解，他将资本划分为若干基本类型，其中最基础的是经济资本、文化资本与社会资本，而象征资本则被视为前三种资本在被社会认可和合法化之后所呈现的表现形式。具体而言，经济资本是指能够直接转化为金钱或制度化财产权的资源，是具有可见性和可量化特征的一类资本。文化资本体现为个体在教育、知识、技能、品味和文化能力等方面的积累，可表现为三种形态，一是具身化形态，指内化于个体之中的形态，无法脱离本人而存在，因而受到个人生物寿命的限制，且需要个人亲自投入时间，不能通过二手方式完成积累或转化，如认知方式与学术素养；二是客观化形态，指具有可传递性的物质形态，如书籍、作品等文化产品；三是制度化形态，指得到制度承认的、以学术资格形式表现的形态，如学历与学位。社会资本是指个体通过持久的社会关系网络所获得的现实或潜在资源。象征资本并非一种独立于其他资本之外的资本形态，而是当某种资本在特定社会情境中被广泛认可并赋予合法性时所呈现的象征性权威与声望\cite{bourdieu1986}。不同类型的资本之间可以相互转化，这一过程遵循着类似能量守恒的法则，资本的每一次转换都需要个人投入相应的劳动时间，付出或多或少的努力才能完成。

布迪厄的资本理论不仅为理解社会结构提供了宏观视角，在高等教育研究的实证领域，其作为分析工具的有效性与适用性也得到了广泛验证。不少学者借助布迪厄的资本理论作为分析框架，探究学者已拥有的或通过受教育获得的资本。Zhang等人将导师的学术资本分为学术声誉和社会资源两个维度，研究发现导师的信息技术整合能力和个人魅力是影响学生感知支持、进而激发创造力的关键\cite{zhang2022}。Walker与Yoon引入"博士资本"的概念，沿用布迪厄的框架，将其拆解为经济资本、社会资本、文化资本和象征资本，指代一组与学术成功相关的各类资本的复合体\cite{walker2017}。Zheng等人在分析中国博士生在东南亚的求学经历时，在经典四重资本的基础上拓展了"语言资本"维度，进一步丰富了该理论的应用\cite{zheng2025}。上述的实证研究表明，布迪厄的资本维度适用于探究学术群体的成长路径，分析学术资本在不同社会与制度情境中的生成、转化机制，因此，本研究在编码框架的构建中，沿用布迪厄对资本形态的经典划分，将CSC资助留学博士生的学术资本划分为社会资本、文化资本和象征资本，以此作为后续质性编码与分析的理论依据。

\subsection{社会资本理论}

社会资本理论（Social Capital Theory）是社会学、政治学与经济学等领域广泛关注的重要理论视角，主要用于解释个体或群体如何通过社会关系网络获取资源，并由此影响其行为选择与发展结果。20世纪后半叶以来，关于社会资本的概念和理论不断丰富和完善，布迪厄首次对社会资本的概念作出了系统阐释，认为社会资本是嵌入在关系网络中、需通过持续投资来维持的实际或潜在资源，并强调社会资本与经济资本、文化资本之间的可转化性。在布迪厄的资本理论框架中，社会资本并非孤立静止的个体所有物，而是通过社会关系的持续积累与维护所生成的关系性资源；它为个体带来象征认可、机会获取与地位巩固等优势。此后，学界在该领域持续深耕，逐步拓展并完善了相关理论体系。科尔曼（Coleman）从功能主义视角出发，将社会资本理解为一种存在于社会结构之中的资源，能够帮助个体实现特定目标，认为闭合的网络比开放的网络更能产生社会资本\cite{coleman1990}。帕特南（Putnam）则从公民参与的角度，将社会资本引入宏观层面，认为社会资本是指社会组织的特征，例如信任、规范和网络，它们能够通过促进协调的行动来提高社会的效率\cite{putnam2000}。林南（Nan Lin）构建了以行动者为中心的分析框架，将社会资本视为个体在有目的的行动中，通过社会网络摄取并使用的资源\cite{linnan2020}。他将社会资本的获取与运用概括为社会关系投资、社会资源的获取与动员，以及由此产生的回报三个相互关联的阶段；从个体视角而言，这一过程也体现为行动者嵌入社会网络、动员社会资源，并最终实现个体行动目标的过程。

社会关系在资源获取与动员中的作用并非同质，围绕这一问题，格兰诺维特（Granovetter）从社会网络结构视角提出了弱关系理论\cite{granovetter1973}。他依据互动频率、情感强度、亲密程度及互惠交换，将社会关系划分为强关系与弱关系，并发现个体间关系越强，其友谊网络重叠度越高，越倾向于形成封闭的小圈子，信息在内部循环；弱关系的力量在于，所有桥梁几乎都是弱关系，弱关系的低重叠度使其能连接不同群体，因此能够在信息传播、社会流动、社区凝聚力中发挥关键作用。后续研究基于格兰诺维特的弱关系理论，通常将社会关系区分为强关系与弱关系，来探究不同强度关系对社会资本获取与使用的影响。强关系连接的个体之间通常互动频繁、情感深厚、信任程度较高，关系长期稳定，并且具有较高的同质性；而形成弱关系的个体之间互动频率较低、情感联结淡、亲密程度较弱，背景差异也相对较大。不少学者将强弱关系理论引入科研领域，既探究两类关系对学者学术产出的差异化影响，也关注学者动态运用强弱关系的策略选择。Abbasi等人证实了格兰诺维特的弱关系理论，发现科学成就与科学家合作网络结构显著相关，拥有弱关系的学者h指数平均更高，且弱关系连接的团队更能产出高被引论文\cite{fronczak2022}。周明睿等人基于深度访谈发现，中国理工科学者一方面依托强关系实施"空间延续"策略以维持既有合作根基，另一方面借助弱关系开展"空间开拓"策略以拓展替代性合作渠道，主动应对地缘政治与疫情叠加下的国际科研合作风险\cite{zhoumingrui2024}。赵鑫与朱佳妮使用混合研究方法，揭示了科研人员在从学者成长为院士的过程中，由师承强关系主导到跨越结构洞的弱关系拓展，再到通过本土学生延续海外强关系的社会资本演化\cite{zhaoxin2025}。

结合研究问题与资料特征，本研究针对性选取格兰诺维特的理论作为分析基础。格兰诺维特提出的强关系与弱关系区分，可明确不同社会关系在资源获取与机会生成中的功能差异，为识别科研社交网络中不同层级的互动提供了分类标准。因此，在编码阶段，本研究以强关系与弱关系理论作为社会关系类型划分的理论依据，对访谈材料中涉及的导师关系、合作关系及学术联系等，按照互动频率、情感强度及亲密度进行类型归纳，确保分类的客观性；在分析阶段，在已分类的基础上，进一步探讨这些不同类型的关系如何具体协助个体积累学术资本，并影响其职业晋升与发展。

\subsection{心理资本理论}

心理资本理论（Psychological Capital Theory）是21世纪初由组织行为学家路桑斯（Fred Luthans）等人基于积极心理学和积极组织行为学（Positive Organizational Behavior，POB）提出的前沿理论，该理论从发展的角度出发，强调个体心理资源的可开发性与增值性，主张通过针对性干预提升个体心理资本水平，进而实现个体与组织的共赢发展。Luthans等人将心理资本定义为个人具有的一种积极的心理发展状态，主要包括自我效能（Efficacy）、乐观（Optimism）、希望（Hope）和韧性（Resilience）四个维度\cite{luthans2018}。具体而言，自我效能概念来源于阿尔伯特·班杜拉（Albert Bandura）的社会认知理论，是指个人对自身能够调动资源、利用所拥有的能力或技能去完成特定任务、达成既定目标的信念与信心。乐观是指个体对于未来发生积极事情的心理预期，并且当事情发生时，主动将积极结果归因于内在、持久性、普遍性的因素，将消极结果归因于外在、暂时性、与情境相关的因素的心理倾向。希望概念来源于Snyder等人提出的"希望理论"（Hope Theory），是一种基于内在成功感的积极动机状态，指个体基于现实且富有挑战性的目标，具备坚定动机与持久意愿，以及能够规划并调整有效达成路径的认知能力。韧性是个体在面对压力、挫折、冲突甚至积极重大的变革时，能够调整心理状态、恢复行动能力，超越平凡，实现正向适应与成长的动态心理能力。

尽管心理资本理论最初应用于企业组织管理领域，但其强调的个体心理资源构建与抗压能力，对于处于跨文化情境下的国际博士生群体同样具有较强的解释力。博士阶段的学术训练过程本身具有高度不确定性、长期投入与多重压力并存的特征，而在国际学术流动背景下，留学海外的博士生不仅要承受学术研究的压力、适应更换导师后的科研与指导模式，还需应对地理迁移与文化差异带来的额外心理挑战，其心理状态及内在调节资源对其学业表现与未来发展具有重要影响。近年来，相关研究逐渐将心理资本理论引入高等教育与研究生培养领域，关注博士生在学术压力情境下的心理资源及其对学术表现的影响。例如，杜学敏、李林英运用扎根理论方法分析优秀博士生的科研实践，指出雄厚的心理资本对于博士生的学术创新具有正向作用\cite{duxuemin2023}。Zhang等人提出了"学业浮力"（Academic Buoyancy）概念，将其视为心理资本在应对日常学业挫折中的具体表现，并发现其在培养学生创造力方面的作用甚至优于社会资本\cite{zhang2024a}。Khuram等人则通过实证研究发现，导师的支持性督导能有效帮助学生积累心理资本，进而提升科研产出\cite{khuram2023}。

因此，本研究将心理资本理论作为一种解释性分析框架，在对国际流动博士生访谈资料进行开放性编码与主轴编码的过程中，关注受访者在跨文化学术流动经历中所形成的积极心理体验与内在调节资源，并在持续比较分析的基础上，对其加以概括和提炼，逐步形成与心理资本内涵相契合的分析范畴。在随后的理论解释阶段，心理资本理论为理解不同积极心理状态如何在国际学术流动情境中生成、累积并发挥作用提供了重要的分析视角，有助于揭示国际流动如何通过影响博士生的心理资本，作用于其学术资本积累与职业发展。

\section{文献综述}

本研究在梳理既有文献时，以布迪厄的资本理论为参照框架，将国际学术流动所产生的影响划分为学术资源及能力积累、科学生产力和职业发展三个方面。在学术资源及能力积累方面，分别梳理了国际流动对学术网络拓展、知识积累与学术声誉三个维度的影响；在科学产出方面，关注国际流动对科研产出数量与质量的影响以及这些影响的条件和限制；在职业发展方面，则聚焦于国际流动经历如何作用于科研人员的职称晋升、机构流动与基金获取等职业结果。本研究在梳理与述评现有研究成果的基础上，结合当前研究的局限开展后续实验设计。

\subsection{国际流动与学术资源及能力积累}

国际学术流动对科研人员学术资源及能力积累的影响，集中体现在学术网络拓展、知识与能力积累以及学术声誉建构三个维度。

在学术网络拓展方面，多数研究从合作者数量、网络结构位置和合作产出等不同角度出发，得出了较为一致的积极结论。Netz等人通过对96项实证研究的系统综述发现，约85\%的研究证实国际流动能够显著促进科研人员国际学术网络的拓展\cite{netz2020}。在因果推断层面，Gu等人运用PSM-DID方法证实CSC资助的国际流动能够增加合作者数量并提升合作产出\cite{gu2024}，Xie与Yegros的统计分析也表明CSC资助学者的国际合作率更高且倾向于与高影响力研究人员合作\cite{xie2025}。从网络结构视角出发，Momeni等人发现学术流动帮助学者占据网络更中心的位置\cite{momeni2022}，Ito等人则揭示了国际流动科学家在连接本土与海外学者中的"桥梁"中介角色\cite{ito2025}。然而，流动效应并非一致正向。Lu与McInerney指出，桥梁位置的优势会随职业阶段变化而递减\cite{lu2016}；赵鑫与朱佳妮发现不同流动目的地催生了不同的网络模式\cite{zhaoxin2025}。在影响因素方面，流动时长、合作持续性和地缘政治等是重要的调节变量------Pan等人发现长期流动（超过12个月）在拓展合作网络方面显著优于短期流动\cite{pan2025}，而Bu等人指出持续合作并非越久越好，其效果因学科性质和团队规模而异\cite{bu2018}。此外，Bauder\cite{bauder2020}、Wang等人\cite{wang2019}和李峰\cite{lifeng2018}的研究均指出，国际流动在拓展新网络的同时可能伴随母国学术网络的流失，Xie与Yegros\cite{xie2025}以及周明睿等人则发现地缘政治变化迫使科研人员调整合作策略\cite{zhoumingrui2024}。从方法论来看，该领域研究以大规模文献计量分析和合著网络分析为主，部分研究采用因果推断方法（如PSM-DID），也有学者运用半结构化访谈探究网络变化的主观机制。

在知识与能力积累方面，相关研究经历了从关注显性产出到重视内在成长的转向。在知识类型化方面，Coey基于对欧洲社会科学与人文学科博士毕业生的访谈，识别出国际流动产生的三种知识成果------显性知识、隐性知识实践与世界主义身份发展，揭示了隐性知识必须深入海外学术环境方能获得的特殊性质\cite{coey2018}。两项综述研究为上述判断提供了系统支撑：Netz等人对96项实证研究的分析表明国际流动能够提升外语水平、技术方法与跨文化协作能力\cite{netz2020}；Shen等人对《Higher Education》期刊的系统分析进一步证实了显性知识与隐性知识的双重积累效应\cite{shen2022}。国内学者则多基于特定院校的数据展开探究，孙伟与赵世奎\cite{sunwei2017}以及姜华等人\cite{jianghua2023}的研究均发现，博士生能够通过海外访学吸收前沿知识、提升科研方法与技能、拓宽学术视野。然而，李澄锋等人运用线性回归分析发现，CSC资助留学博士生与未参与留学的博士生在科研能力增值方面并无显著差异\cite{lichengfeng2020}，提示国际流动对能力发展的作用可能受到流动时长、嵌入深度和学科特性等因素的调节。从方法论来看，该领域以定性访谈和系统综述为主，定量研究多采用问卷调查，科研能力主要通过主观自我评价测度，在具体能力形成过程与不同情境下的差异机制方面仍缺乏更为细致的经验解释。

在学术声誉方面，现有研究主要从制度性认可和学术产出质量两条路径展开分析。在制度性认可路径上，学者们较为一致地发现国际流动经历本身即具有声誉价值：Bauder等人从批判性视角提出"流动崇拜"（mobility fetishism）概念，指出学术评价体系中国际流动经验已被视为衡量学术卓越性的重要指标\cite{bauder2017}；Tran\cite{tran2016}和Leung\cite{leung2013}分别通过对澳大利亚国际学生和赴德中国学者的访谈，证实海外经历被视为"区分标志"，其声誉价值随机构声望的提升而增强。在学术产出质量路径上，Biazoli等人\cite{biazoli2025}和Li等人\cite{li2025}的研究分别表明，国际流动能够提升发表期刊的学术声望、论文引用量以及在高质量期刊独立发表的可能性。然而，这一积极效应并不绝对------一项针对爱尔兰学者的研究指出，学术声誉的生成并非取决于流动经历本身，而是流动经历与接收场域互动的结果，只有在海外经历被认可为高质量科研经历时，才能真正转化为学术声誉\cite{courtois2024}。从方法论来看，该领域以定性访谈为主，定量研究因难以直接测量学术声誉，多以期刊声望、引用量等替代变量进行间接检验，对声誉形成的过程机制探讨仍较为有限。

\subsection{国际流动与科学生产力}

围绕国际学术流动与科研产出的关系，学界积累了相当数量的实证研究。多项研究从不同国家情境出发，得出了较为一致的积极结论：Xie与Yegros基于文献计量分析发现CSC资助者的引用影响力高于中国及全球平均水平\cite{xie2025}，Gu等人运用PSM-DID方法证实CSC资助的国际流动能显著提升研究产出数量与质量并产生持续约4年的积极影响\cite{gu2024}，Aykac对土耳其科研人员的因果推断分析\cite{aykac2021}以及Miranda与Cabello对巴西博士生导师的研究\cite{grochocki2022}也得出了类似结论。然而，这一积极效应并非普遍存在。叶晓梅和梁文艳的研究表明海归教师与本土教师在科研产出方面差距不大\cite{yexiaomei2019}，Chu等人则发现在硬学科中本土博士的科研生产力甚至高于留学归国博士\cite{chu2024}。值得注意的是，李澄锋等人\cite{lichengfeng2020}和Bao等人\cite{bao2023}的研究揭示了产出结构的差异------国际流动主要提升国际期刊发表，而对国内期刊发表的促进作用有限，表明国际流动更多改变了学者的发表模式而非总体产出规模。

在影响因素方面，流动时长、流动目的地、职业阶段和学科背景等是调节国际流动效应的关键变量。Gao与Liu的访谈研究揭示了流动时长的双面性------长期海外经历有助于英语写作优势的形成，而短期访学的效果因职业阶段而异\cite{gao2021}。Han等人基于顶尖学者的大规模数据发现流动目的地的效应存在显著差异，以美国的提升效果最为明显\cite{han2024}。李峰进一步指出国际流动对科研产出具有阶段性特征，呈现"短期回落、长期提升"的时间模式\cite{lifeng2018}。张颖等人则发现流动空间跨度对国际合作产出存在一定抑制作用\cite{zhangying2024}。上述研究表明，国际流动与科研产出之间并非简单的线性关系，其效果受到多重因素的共同调节。

从方法论来看，该领域以定量研究为主导，广泛采用文献计量分析、回归分析和因果推断方法（如PSM-DID、DID），以论文数量、引用量、期刊影响因子等客观指标测度科研生产力。也有少数学者采用访谈法探究产出差异背后的机制。相较于学术网络和知识积累领域，科学生产力研究在方法的严谨性和可比性方面更具优势，但对于产出背后的微观过程------如国际流动如何通过知识积累与网络拓展间接促进科研产出------关注相对不足。

\subsection{国际流动与职业发展}

围绕国际学术流动对职业发展的影响，学界形成了积极效应与消极效应并存的研究格局。从积极效应来看，多项研究从职称晋升、初职获取和基金资助等不同角度出发，得出了较为一致的正向结论。Zweig等人发现海归博士的职业晋升速度快于本土博士，且在申请和获得国家级科研项目资助方面更具优势\cite{zweig2004}；黄小燕等人对经济学与管理学领域博士生的研究表明，CSC资助的博士生项目有助于科研人员初职进入高声望机构、加快早期职称晋升并提升获得国家级科学基金资助的可能性\cite{huangxiaoyan2026}；施云燕基于问卷调研发现国际流动更有助于学术职业发展而非行政领域晋升，且随着职称等级的提升，认为国际流动具有积极作用的科研人员占比呈上升趋势\cite{shiyunyan2017}。Lu与McInerney则从网络结构视角出发，发现归国学者可借助结构洞位置获得职业优势，且在终身教职体系中国际学术网络的促进作用尤为突出\cite{lu2016}。然而，另一些研究得出了截然不同甚至相反的结论。Van Mol等人考察荷兰高等教育毕业生后发现，尽管描述性统计显示流动学生找到第一份工作的时间更短，但在控制选择效应后这一优势完全消失，说明积极效应可能源于群体自选择而非流动本身\cite{vanmol2021}。Cruz-Castro等人基于对西班牙科学家的研究发现，在高度强调机构忠诚的制度体系下，国际流动并未对科研人员获得早期终身教职产生积极作用\cite{cruzcastro2010}。Seeber等人针对比利时弗拉芒地区大学的定量分析同样显示，国际流动会降低科研人员在原机构获得晋升的可能性，同时提高其离职概率\cite{seeber2023}。上述分歧表明，国际流动对职业发展的影响并不是一种普遍性的正向效应，而是高度依赖于制度情境与个体条件。

在影响因素方面，制度环境、流动类型和学科差异是调节国际流动职业效应的三个关键维度。在制度环境方面，Bauder通过对加拿大和德国国际流动学术人员的半结构化访谈发现，在加拿大，国际流动通常促进职业发展，而在德国，国际流动常导致母国网络流失并对职业发展产生负面影响，揭示了国家层面学术劳动力市场结构与制度认可机制对流动效应的根本性制约\cite{bauder2020}。在流动类型方面，潘虹与唐莉针对长江学者特聘教授的研究表明，海外攻读博士及境外任教经历均有助于推动职业发展，而短期访学或长期在境外工作则不利于学术资源的积累，反而阻碍职业发展\cite{panhong2023}；李峰与唐莉的研究进一步发现，全职海外工作经历对学者职业晋升具有显著加速效应，尤其是在海外取得初级职称后归国的学者优势最为突出，而仅拥有海外博士学位或短期访学经历者对职业发展速度并无显著促进\cite{lifeng2022}。在学科差异方面，王艳与赵江南基于Nature全球博士后调查数据的研究表明，国际流动在医学领域对职业发展呈现负向效应，在理工科影响有限，而在社会科学领域则表现出积极的促进作用，这一发现提示不同学科的知识生产模式和职业评价标准会显著调节国际流动的职业回报\cite{wangyan2026}。

从方法论来看，该领域研究呈现定量与定性方法并重的特点。定量研究多采用问卷调查和统计分析，以职称等级、晋升时间、基金获取等客观指标测度职业发展结果，部分研究尝试通过控制选择效应来增强因果推断的有效性。定性研究则主要采用半结构化访谈，侧重于揭示职业发展的过程性机制与制度情境的交互作用。然而，现有研究大多从横截面视角考察流动与职业结果的关系，对于国际流动如何通过学术资源积累这一中间环节渐进地影响职业发展的纵向过程，实证追踪仍较为匮乏。

\subsection{文献述评}

综合上述三个方面的文献梳理可以看出，学界围绕国际学术流动效应已积累了大量实证证据，研究主题覆盖了从微观层面的个体资源获取到宏观层面的职业结果变化，研究方法也从早期以描述性分析为主逐步发展为运用因果推断、文献计量和深度访谈等多元手段的综合研究格局。多数研究证实国际流动能够在不同维度上为科研人员带来积极影响，但也有相当数量的研究揭示了流动效应的条件性、异质性乃至负面可能。尽管如此，现有研究仍存在若干有待进一步深化之处。

首先，现有研究尚未从资本积累的角度对科研人员国际学术流动效应进行系统研究。既有文献较多采用量化方法考察国际流动与科研产出、合作关系或职业结果之间的相关性，但对于科研人员在国际流动过程中究竟积累了哪些类型的学术资本，以及不同类型资本如何在流动过程中生成与转化，缺乏系统性的理论框架与经验分析。尤其是在隐性知识如何传递与吸收、学术网络如何由"建立"转向"有效利用"等方面，现有研究仍缺乏更为细致的分析。此外，心理资本维度在相关研究中尚未得到充分重视。与学术网络、知识积累和学术声誉等方面相比，国际流动对科研人员乐观、希望、自我效能感和韧性等心理资源的影响讨论相对较少。国际流动经历是否会促进心理资本积累，心理资本又如何与其他方面的积累相互作用并进一步影响科研人员的学术成长与职业发展，这些问题仍有待深入探讨。

其次，现有研究对早期学术积累如何影响后续职业发展的过程性机制关注不足。国内现有研究主要以已经获得一定成就的科研人员为研究对象，如长江学者、杰青获得者等，较少关注成长型科研人员的国际流动经历及其影响。然而，处于职业早期阶段的科研人员正是国际学术流动的主要参与者，其在流动过程中积累的学术资源与能力如何转化为后续职业发展的竞争优势，这一从积累到转化的过程性机制仍有待深入揭示。

最后，现有研究较少关注科研人员博士阶段的国际流动带来的长期效应。博士阶段是科研人员学术训练与职业准备的关键时期，也是国际学术流动最为活跃的阶段之一，但既有研究多将不同职业阶段的流动经历混合讨论，较少单独考察博士阶段国际流动对科研人员后续学术成长与职业发展的持续性影响。关于国家留学基金委资助博士生的研究虽已逐步展开，但多数研究仍集中于特定学科或个别院校，缺乏针对该群体的系统性和整合性考察。

基于上述研究现状与不足，本文拟从以下几个方面对现有研究作出补充。首先，在理论框架上，本文将在布迪厄资本理论基础上，结合心理资本相关研究，从文化资本、社会资本、象征资本和心理资本四个维度构建分析框架，以考察国际学术流动对科研人员早期学术资本积累的影响。其次，在研究方法上，本文采用质性研究路径，通过深度访谈还原流动者的具体经历，从微观层面分析"国际流动---学术资本积累---职业发展"的过程机制，探讨资本积累的形成路径、转化条件及其阻碍因素。最后，在研究对象上，本文聚焦于国家留学基金委资助的博士研究生，系统考察其博士阶段国际流动经历对早期学术资本积累与职业发展的长期影响，以期为理解我国公派留学背景下成长型科研人员的发展过程提供经验材料，并为相关政策完善和博士生培养实践优化提供参考。
